\chapter{Integrazione Hub futura}
\label{ch:hub}

\section{Endpoint}
Nuova Hub obbligatoria (la vecchia polyglot risponde 530):
console \url{https://gamification-webapp.createlab-univaq.it},
API \url{https://gamification-api.createlab-univaq.it/api/v1},
Swagger \url{https://gamification-api.createlab-univaq.it/swagger-ui/index.html},
guida \url{https://gamification-api.createlab-univaq.it/docs/api-guide.en.md}.
Auth: \texttt{POST /auth \{username, password, origin:''GAME''\}}
$\rightarrow$ Bearer JWT (24\,h). Runtime: \texttt{POST /executions
\{gameId, playerId, actionId, data:\{\}\}} (il campo \texttt{data} è
obbligatorio anche se vuoto).

\section{Concetti Hub}
Actions, point concepts con periodi in millisecondi, badge come collezioni,
un solo level per point concept, regole Drools con \texttt{/rules/validate}
e simulazioni, challenge / leaderboard / classifiche.
\textbf{Migrazione}: i vecchi \texttt{gameId} non sono validi; il gioco va
ricreato da console e le regole riscritte.

\section{Payload Dart di esempio}
\begin{lstlisting}[caption={Esempio di chiamata executions con package:http.},label={lst:exec}]
import 'package:http/http.dart' as http;
import 'dart:convert';

Future<void> sendAction({
  required String base, required String token,
  required String gameId, required String playerId,
  required String actionId, Map<String, dynamic> data = const {},
}) async {
  final res = await http.post(
    Uri.parse('$base/api/v1/executions'),
    headers: {'Content-Type': 'application/json',
              'Authorization': 'Bearer $token'},
    body: jsonEncode({'gameId': gameId, 'playerId': playerId,
                      'actionId': actionId, 'data': data}),
  );
  if (res.statusCode != 200 && res.statusCode != 201) {
    throw StateError('Hub error ${res.statusCode}: ${res.body}');
  }
}
\end{lstlisting}

\section{Strategia di migrazione}
\texttt{FakeEngineClient} che implementa l'interfaccia dei repository e
instrada ogni azione locale verso un \texttt{actionId} Hub
(\texttt{complete\_quiz} $\rightarrow$ punti XP, \texttt{unlock\_topic}
$\rightarrow$ badge/prerequisiti, \texttt{boss\_hit} $\rightarrow$ punti
boss, \texttt{earn\_card} $\rightarrow$ collezione). I mock restano come
fallback offline. Mapping \texttt{actionId} definitivo da fissare in console
al momento della creazione del gioco.
